 boundary, while every individual logit-score
norm remains bounded.
\end{theorem}

The proof is provided in
Appendix~\ref{app:policy_family_manifestations}.
Fixed covariance is sufficient for the Gaussian result; learned covariance is
not required. Neither policy-family conclusion alone implies
task-performance collapse or a NaN/Inf numerical failure.

\subsection{Why Selective Tapering Works}
\label{sec:selective_tapering_bridge}

The preceding regimes identify the failure; the following radial limit
identifies the control property that reverses it.

\begin{lemma}[Far-field radial balance]
\label{lem:far_field_radial_balance}
For a fixed-covariance Gaussian field with finite negative atoms,
\[
\begin{aligned}
F(\mu)&=p\Sigma^{-1}(m_+-\mu)
 +\sum_iq_i\Sigma^{-1}(\mu-a_i),\\
q&=\sum_iq_i,
\end{aligned}
\]
and \(\mu=m_++rv\), \(\|v\|_2=1\),
\[
\lim_{r\to\infty}\frac{\langle v,F(\mu)\rangle}{r}
=(q-p)v^\top\Sigma^{-1}v.
\]
\end{lemma}

\begin{corollary}[Selective tapering versus global scaling]
\label{cor:selective_vs_global}
If each negative atom is multiplied by
\(\omega_i(D_i)\to0\), the limit in
Lemma~\ref{lem:far_field_radial_balance} becomes
\(-p v^\top\Sigma^{-1}v<0\). Thus, for every fixed finite
\(\rho=q/p\), without \(\rho\)-specific tuning, the far-field radial
component is eventually inward; the onset radius is not uniform in
\(\rho\) and depends on the tail rate of \(\omega\). If additionally
\(\omega(D)\sqrt D\to0\), the absolute tapered negative force vanishes.
In contrast, global scaling \(\omega\equiv\alpha\) is inward only when
\(\alpha q<p\). Proofs are in Appendix~\ref{app:far_field_bridge}.
\end{corollary}

\begin{proposition}[Restoration by reference regularization]
\label{prop:regularized_restoration}
The regularized linear field has a unique finite stable equilibrium when
\(p>q\) for every \(\gamma\ge0\), when \(p=q\) for every \(\gamma>0\),
and when \(p<q\) if and only if \(\gamma>q-p\). At \(p=q\), its
displacement is \(O(\gamma^{-1})\). Reference regularization restores
stability by tethering the policy; DRPO instead removes the far-field
source. The proof and the Euclidean-spring variant are in
Appendix~\ref{app:regularized_restoration}.
\end{proposition}

Corollary~\ref{cor:selective_vs_global} provides the mechanism-level
justification for Section~\ref{sec:drpo}: selective tapering changes the
far-field sign by removing the remote source, whereas a constant global scale
must be tuned to the negative-to-positive mass ratio. These conclusions assume
frozen data, fixed advantages, and unbounded reuse. Data refresh, actor--critic
coupling, and reference regularization change the field; the proposition above
shows how regularization instead restores stability by tethering the policy.

\section{Dynamic Remoteness-Aware Policy Optimization}
\label{sec:drpo}

Section~\ref{sec:selective_tapering_bridge} identifies the design target:
control the remote negative source without uniformly weakening the negative
branch. We first define DRPO, then derive its control-order hierarchy and
closed-loop guarantees.

\subsection{Update and Policy-Family Instantiations}
\label{subsec:drpo_update}

For remoteness \(D\), DRPO uses the thresholded exponential taper
\begin{equation}
    \omega_{\mathrm{Exp}}(D)
    =
    \exp
    \left\{
        -\lambda
        \left[
            \frac{D-\tau}{c}
        \right]_{+}
    \right\},
    \label{eq:drpo_exp_weight}
\end{equation}
where \(\tau\) marks the control threshold, \(c>0\) fixes the remoteness
scale, and \(\lambda>0\) controls the decay rate. The weight is exactly one
when \(D\leq\tau\) and decays smoothly when \(D>\tau\).

Using \(D_\theta(s,a)=-\log\pi_\theta(a\mid s)\), define
\(\bar D_\theta(s,a)=\operatorname{sg}[D_\theta(s,a)]\), where
\(\operatorname{sg}[\cdot]\) denotes stop-gradient. DRPO changes only the
negative branch:
\begin{equation}
    \widetilde A_{\theta}(s,a)
    =
    \widehat A^{+}(s,a)
    -
    \omega_{\mathrm{Exp}}
    \!\left(\bar D_\theta(s,a)\right)
    \widehat A^{-}(s,a),
    \label{eq:drpo_effective_advantage}
\end{equation}
with actor field
\begin{equation}
    \mathbf F_{\mathrm{DRPO}}(\theta)
    =
    \mathbb E_{(s,a)\sim\nu}
    \left[
        \widetilde A_{\theta}(s,a)
        \nabla_\theta\log\pi_\theta(a\mid s)
    \right].
    \label{eq:drpo_actor_field}
\end{equation}
Remoteness is recomputed under the current policy but detached before it is
used as a coefficient; DRPO is therefore a dynamic weighting rule rather than
a differentiable distance regularizer.

At each policy iterate, tapering preserves the algebraic regime structure of
Theorem~\ref{thm:aggregate_equilibria} after replacing the negative mass and
moment by
\[
q_\omega=\mathbb E_\nu[\widehat A^-\omega],
\qquad
\mathbf M_{-,\omega}=\mathbb E_\nu[\widehat A^-\omega T(a)].
\]
When \(0<q_\omega<p\), the controlled signed target is
\(\mathbf m_\omega^\star=(p\mathbf m_+-\mathbf M_{-,\omega})/(p-q_\omega)\);
when \(q_\omega=0\), use \(\mathbf M_{-,\omega}=0\) without defining a
normalized negative moment.

\paragraph{Gaussian policies.}
For a Gaussian policy with mean \(\mu_\theta(s)\) and fixed covariance
\(\Sigma\), define
\begin{equation}
    D_{\mathrm G,\theta}(s,a)
    =
    (a-\mu_\theta(s))^\top
    \Sigma^{-1}
    (a-\mu_\theta(s)).
    \label{eq:gaussian_squared_distance}
\end{equation}
Since \(D_\theta-C_\Sigma=\frac12D_{\mathrm G,\theta}\), the additive
constant and factor \(1/2\) can be absorbed into the threshold and scale:
\begin{equation}
    \omega_{\mathrm G}(s,a)
    =
    \exp
    \left\{
        -\lambda_{\mathrm G}
        \left[
            \frac{D_{\mathrm G,\theta}(s,a)-\tau_{\mathrm G}}
                 {c_{\mathrm G}}
        \right]_{+}
    \right\}.
    \label{eq:drpo_gaussian_weight}
\end{equation}

\paragraph{Categorical policies.}
For a categorical policy, substituting selected-action surprisal into
Equation~\eqref{eq:drpo_exp_weight} gives
\begin{equation}
    \omega_{\mathrm C}(s,a)
    =
    \min
    \left\{
        1,\,
        \left(e^{\tau}\pi_\theta(a\mid s)\right)^{\lambda/c}
    \right\}.
    \label{eq:drpo_categorical_weight}
\end{equation}
Thus, the same actor update uses squared standardized distance for Gaussian
policies and surprisal for categorical policies.

\subsection{Loop-Gain Order}
\label{subsec:min_order_control}

The quantity that drives the next change in remoteness is not the
single-step parameter displacement but the squared score
\(R(D)=\|\nabla D\|_2^2\). We call this the reuse \emph{loop gain}.

\begin{proposition}[Gaussian loop-gain order]
\label{prop:loop_gain_order}
For a fixed-covariance Gaussian mean policy,
\(I_{\mathrm{loop}}(D):=R(D)=\Theta(D)\) in the far field. Hence
\(\omega(D)I_{\mathrm{loop}}(D)\) is bounded exactly at the matching
order \(\omega(D)=O(D^{-1})\), and vanishes when
\(\omega(D)=o(D^{-1})\). Equivalently, the matching order is
\(O(r^{-2})\) in standardized distance \(r\), since
\(D-C_\Sigma=r^2/2\). The proof is in
Appendix~\ref{app:min_order_control}.
\end{proposition}

At the exact matching rate \(\omega(D)=\Theta(D^{-1})\), the weighted
parameter force is \(\Theta(D^{-1/2})\to0\), but the weighted reuse loop gain
is \(\Theta(1)\). An isolated reused atom can therefore keep driving
remoteness outward at a constant order. Quantitatively, uncontrolled Gaussian
reuse gives geometric growth, matching-order tapering gives linear growth,
and DRPO's exponential taper gives logarithmic growth.

\subsection{Distance-Dependent Utility of Negative Feedback}
\label{subsec:distance_dependent_utility}

Bounded loop gain alone does not determine whether a remote update is useful.
Appendix~\ref{app:negative_update_utility} gives an empirical utility
definition and a sufficient shrinkage result, used here only as motivation;
we do not assume a universal monotone utility law in remoteness. When a
historical negative no longer supplies local correction information, the
design target is
\begin{equation}
    \omega(D) I(D)
    \longrightarrow 0
    \qquad
    \text{as }
    D\longrightarrow\infty .
    \label{eq:far_field_residual_vanishes}
\end{equation}
This is a conditional design principle, not a claim that every remote
sample is harmful. If far-field feedback remains aligned with the true
task-improving direction, tapering can remove useful signal.

\subsection{Guarantee Hierarchy}
\label{sec:exponential_taper}

The DRPO taper satisfies the stronger requirement above for every
finite-order response considered here.

\begin{proposition}[Vanishing finite-order far-field response]
\label{prop:exp_vanishing_influence}
Let \(I(D)\) denote the raw far-field response scale. If, for some
finite \(m\geq0\) and \(C>0\),
\(I(D)\leq C(1+D)^m\) in the far field, then the taper in
Equation~\eqref{eq:drpo_exp_weight} satisfies
\[
    \omega_{\mathrm{Exp}}(D) I(D)
    \longrightarrow 0
    \qquad
    \text{as }D\to\infty .
\]
Thus, the exponential taper eliminates any finite-order far-field tail
rather than merely bounding it.
\end{proposition}

The exponential form is not unique; it is a simple smooth envelope that
dominates finite-order growth. The proof is given in
Appendix~\ref{app:exp_vanishing_influence}.

The control conditions form a guarantee hierarchy. L0,
\(\omega(D)\to0\), gives eventual inward Gaussian radial drift for each
fixed finite \(\rho\); L1, \(\omega(D)\sqrt D\to0\), removes the
absolute negative far-field force; L2, \(\omega(D)D=O(1)\), bounds the
Gaussian reuse loop gain; and L3,
\(D^k\omega(D)\to0\) for every finite \(k\), removes all finite-order
tails. Global scaling fails L0, reciprocal-quadratic tapering is critical at
L2, and exponential tapering satisfies L0--L3. This hierarchy ranks
guarantees, not task performance.

This hierarchy also organizes the experimental controls. Positive-only and
global scaling ignore remoteness; reciprocal tapers have polynomial tails, and the quadratic member meets the matching
order; DRPO uses the exponential member that satisfies
L0--L3. These controls share the actor form in
Equation~\eqref{eq:drpo_actor_field} and differ only in their negative-branch
coefficient \citep{peng2019advantage,nair2020awac,kostrikov2021offline,
grigsby2021closer,arnal2025asymmetric,leroux2025tapered}.

\subsection{Output-Coordinate Closed-Loop Guarantees}

We now record the output-coordinate consequences of the DRPO update once
positive attraction is included.

\begin{theorem}[Gaussian DRPO boundedness]
\label{thm:gaussian_drpo_boundedness}
Let \(p>0\), \(\Sigma\succ0\), and
\[
F_\omega(\mu)=p\Sigma^{-1}(m_+-\mu)
+\sum_iq_i\omega_i(D_\mu(a_i))\Sigma^{-1}(\mu-a_i).
\]
For finitely many negative atoms with exponential tapers, define
\begin{equation}
\begin{aligned}
B_\omega&=\sup_\mu\left\|\sum_iq_i\omega_i(D_i)
\Sigma^{-1}(\mu-a_i)\right\|_2,\\
R^*&=\frac{B_\omega}{p\lambda_{\min}(\Sigma^{-1})}.
\end{aligned}
\label{eq:gaussian_bound_radius}
\end{equation}
Every continuous trajectory satisfies
\(\limsup_{t\to\infty}\|\mu(t)-m_+\|_2\le R^*\), and the field has at
least one finite equilibrium. If an uncontrolled stable equilibrium
places every negative atom strictly below its threshold, the controlled
field and Jacobian coincide with the original ones in a neighborhood of
that equilibrium.
\end{theorem}

The proof gives an explicit branch-wise bound on \(B_\omega\); a separate
branch is required when the shifted threshold is negative because the taper is
then active at radius zero. Along every escaping sequence,
\(\rho_{\mathrm{eff}}(\mu)=p^{-1}\sum_iq_i\omega_i(D_i)\to0\). This is a
diagnostic consequence, not the boundedness argument. See
Appendix~\ref{app:gaussian_drpo_boundedness}.

\begin{theorem}[Categorical self-limiting suppression]
\label{thm:categorical_self_limiting}
For a repeatedly suppressed action \(i\), define
\(y_i=z_i-\log\sum_{j\ne i}e^{z_j}\), so \(\pi_i=\sigma(y_i)\).
Under an uncontrolled negative atom of mass \(q>0\),
\[
\dot y_i=-q(1-\pi_i)c_K(t),
\qquad \frac K{K-1}\le c_K(t)\le2.
\]
Hence \(y_i\) eventually decreases linearly and \(\pi_i\) exponentially.
If one off-target action is initially the unique maximum then
\(c_K(t)\to2\); under complete off-target symmetry,
\(c_K(t)\equiv K/(K-1)\). With
\(\omega_i=\min\{1,(e^{\tau_i}\pi_i)^\beta\}\), \(\beta=\lambda/c\),
the far branch instead satisfies
\[
y_i(t)=-\beta^{-1}\log t+O(1),
\qquad \pi_i(t)=\Theta(t^{-1/\beta}).
\]
Thus tapering changes persistent suppression from exponential to
polynomial probability decay; it does not freeze the action probability.
The output-coordinate proof is in
Appendix~\ref{app:categorical_self_limiting}; shared-network parameter
gradients additionally contain the network Jacobian and cross-sample
interference.
\end{theorem}

\begin{remark}[Restoring-force balance]
A stable finite crossing requires the restoring-minus-suppression field to
change sign from positive to negative as the log-odds increase; equivalently,
at a hyperbolic crossing, \([S_{\mathrm{tap}}-E]'(y^*)>0\). Entropy regularization with
\(\beta\ge1\) satisfies only the required far-tail ordering in the binary or
symmetric one-vs-rest reduction; a full stability result additionally uses
the stated phase-line conditions. See
Appendix~\ref{app:categorical_restoring_balance}.
\end{remark}

\subsection{Variational Interpretation}

DRPO's exponential attenuation admits an optimistic finite-measure
reweighting interpretation analogous to divergence-ball DRO. The analogy
is deliberately qualified: unlike standard ambiguity sets over normalized
probability measures, the present construction permits attenuation of the
total negative-update mass. Let
\(d\xi^-(z)=\widehat A^-(z)d\nu(z)\) and
\(h(z)=[D(z)-\tau(z)]_+\). Consider
\begin{equation}
\min_{0\le w\le1}\int wh\,d\xi^-
\quad\mathrm{s.t.}\quad
\int(w\log w-w+1)d\xi^-\le\delta.
\label{eq:finite_measure_variational}
\end{equation}

\begin{proposition}[Finite-measure exponential reweighting]
\label{prop:finite_measure_variational}
Let \(\delta_0=\xi^-(\{h>0\})\). For
\(0<\delta<\delta_0\), the constraint in
Equation~\eqref{eq:finite_measure_variational} is active and has a unique
finite \(\kappa>0\) with
\(w_\kappa(z)=e^{-h(z)/\kappa}\). At \(\delta=\delta_0\), the endpoint is
\(w_0=\mathbf1\{h=0\}\); for \(\delta>\delta_0\), zero objective is
already attainable and no finite one-to-one \(\kappa\)--\(\delta\)
relation remains. The later endpoint \(\delta\ge|\xi^-|\) merely makes
\(w\equiv0\) feasible. The proof is in
Appendix~\ref{app:finite_measure_variational}.
\end{proposition}

\subsection{Relation to TOPR}

For sequence policies, token surprisal sums to total trajectory surprisal
\(-\log\pi_\theta(y\mid x)\); mean-token surprisal divides this quantity by
\(|y|\). The two coordinates have the same ordering at fixed length but need
not agree across variable-length responses. TOPR attenuates negative updates
by the truncated ratio \(\min\{\pi/\mu,1\}\).

At the atomic-action or total-trajectory-surprisal level, canonical TOPR's
negative coefficient is the 